% Demo / documentation deck for the modern `illinois` beamer theme.
%
% Compile with:  pdflatex illinois-demo.tex   (run twice for frame totals)
%
% This file doubles as a quick reference for the theme's features.

\documentclass[aspectratio=169]{beamer}

% Load the theme. Options:
%   \usetheme[nowordmark]{illinois}      % hide the title-page wordmark
%   \usetheme[nosectionpages]{illinois}  % no automatic section slides
\usetheme{illinois}

\title{The \texttt{illinois} Beamer Theme}
\subtitle{A modern University of Illinois look}
\author{John and Mary Doe}
% The optional short form (in brackets) is shown in the footline.
\institute[UIUC]{University of Illinois Urbana-Champaign}
\date{\today}

\begin{document}

\frame{\titlepage}

\section{Getting Started}

\begin{frame}{Using the theme}
  Add the theme to any \texttt{beamer} document:
  \begin{block}{Preamble}
    \texttt{\textbackslash documentclass\{beamer\}}\\
    \texttt{\textbackslash usetheme\{illinois\}}
  \end{block}
  It is self-contained --- it needs only \texttt{beamer}, \texttt{xcolor},
  \texttt{graphicx}, and \texttt{hyperref}.
\end{frame}

\begin{frame}{Lists}
  \textbf{Unordered}
  \begin{itemize}
    \item University of Illinois Urbana-Champaign
    \item Department of Statistics
    \item Illinois Informatics Institute
  \end{itemize}

  \textbf{Ordered}
  \begin{enumerate}
    \item First point
    \item Second point
  \end{enumerate}
\end{frame}

\section{Blocks and Math}

\begin{frame}{Blocks}
  \begin{block}{Standard block}
    Illinois blue title, soft tinted body.
  \end{block}
  \begin{alertblock}{Alert block}
    Illinois orange for emphasis.
  \end{alertblock}
  \begin{exampleblock}{Example block}
    Used for examples and asides.
  \end{exampleblock}
\end{frame}

\begin{frame}{Mathematics}
  \begin{exampleblock}{Binomial Theorem}
    \begin{equation*}
      f\left(k\right) = \binom{n}{k} p^{k}\left(1-p\right)^{n-k}
    \end{equation*}
  \end{exampleblock}
  Inline math such as $e^{i\pi} + 1 = 0$ is colored normally.
\end{frame}

\begin{frame}{Columns}
  \begin{columns}
    \begin{column}{0.48\textwidth}
      \textit{Hello!}
      \begin{itemize}
        \item left column
      \end{itemize}
    \end{column}
    \begin{column}{0.48\textwidth}
      \textbf{Goodbye!}
      \begin{enumerate}
        \item right column
      \end{enumerate}
    \end{column}
  \end{columns}
\end{frame}

\end{document}
